\documentclass[10pt,a4paper]{amsart}
\usepackage[margin=19mm]{geometry}
\usepackage{amsmath,amssymb,mathtools}
\usepackage[scr=rsfs,cal=euler]{mathalfa}
\usepackage{xfrac}
\usepackage[unicode,hidelinks,bookmarks=false]{hyperref}
\newcommand{\ie}{i.e.\ }
\newcommand{\cf}{cf.\ }
\newcommand{\resp}{respectively\ }
\newcommand{\Subsection}{\subsection}
\providecommand{\nobreakdash}{\nobreak\hbox{-}}
\setlength{\emergencystretch}{2em}
\multlinegap0pt
\usepackage{amssymb}
\newtheorem{conj}{Conjecture}
\newtheorem{lem}{Lemma}
\newtheorem{prop}{Proposition}
\newtheorem{thm}{Theorem}
\newcommand\F{\mathbb{F}}
\newcommand\Symgrp{\operatorname{Sym}}
\newcommand\diam{\operatorname{diam}}
\title{Common neighbour conjectures for Saxl graphs fail at every base size}
\author{Aluna Rizzoli, Adam R. Thomas}
\date{Source: arXiv:2609.01367v1 (CC BY 4.0)}
\begin{document}
\maketitle

\noindent\emph{Source and licence.} Common neighbour conjectures for Saxl graphs fail at every base size. Version arXiv:2609.01367v1 (CC BY 4.0), distributed by arXiv under the Creative Commons Attribution 4.0 International licence, http://creativecommons.org/licenses/by/4.0/, as recorded in the arXiv licence metadata for this exact version and shown on the abstract page https://arxiv.org/abs/2609.01367v1. Full licence notice: \texttt{licenses/arxiv-2609.01367-CC-BY-4.0.txt}.\par\smallskip

\noindent\textit{Editorial scope.} Counted unit: the article's thm thm:affine (source0.tex lines 264--266). The article's complete original proof of it is delivered with it (source0.tex lines 268--281, one \texttt{proof} environment of 14 lines). No mathematics is written, repaired or re-derived: the statement and the proof are the article's own lines, in the article's own notation, and no retained line is substituted or re-flowed.  Retained uncounted as the source-local context the proof needs: lines 54--56; lines 107--109; lines 118--120; lines 124--129; lines 133--135; lines 143--145; lines 191--199; lines 243--249; lines 252--254; lines 257--261.  Nothing was omitted from any retained span, so the package carries no omission record.  No retained line contains a citation, so the wrapper carries no bibliography and no bibliography entry.  No retained line contains a figure environment, an image-inclusion command or drawing code, so no image is delivered and the package declares no asset dependency.  Editorial formatting only, in addition: an AMS \texttt{amsart} preamble that restates the article's own declarations and macros used by the retained lines, namely the environments \texttt{conj}, \texttt{lem}, \texttt{prop}, \texttt{thm} on separate counters, together with the standard packages those lines need.  The retained environments print in this document as Conjecture 1, Lemma 1, Proposition 1, Theorem 1 in order of appearance, whereas the article prints the same items with the numbering of its own full document.  The complete native source of the article as distributed in the arXiv e-print for this exact version is bundled unchanged as \texttt{sources/209104/source0.tex}.
\begin{conj}[Common Neighbour Conjecture]\label{bgconj}
Let $G$ be a primitive permutation group with $b(G)=2$. Then any two vertices of $\Sigma(G)$ have a common neighbour.
\end{conj}

\begin{lem}\label{paprim}
Let $X\le\Symgrp(\Delta)$ and $Q\le S_m$ be nontrivial, where $m\ge2$. Then $X\wr Q$ is primitive on $\Delta^m$ in product action if and only if $X$ is primitive but not regular on $\Delta$ and $Q$ is transitive on $\{1,\dots,m\}$.
\end{lem}

\begin{lem}\label{basecrit}
Let $x,y\in\Delta^m$. Then $\{x,y\}$ is a base for $G=X\wr Q$ if and only if each $(x_j,y_j)$ is an ordered base for $X$ and the partition $\Pi(\rho_{x,y})$ is distinguishing for $Q$.
\end{lem}

\begin{prop}\label{bookkeeping}
Let $X$ be transitive with $b(X)=2$ and let $Q\le S_m$ be transitive. Then the map induced by $(x,y)\mapsto\rho_{x,y}$ identifies the elements of $\mathcal R(X\wr Q)$ with the $Q$-orbits of tuples $\rho\in\mathcal R(X)^m$ for which $\Pi(\rho)$ is distinguishing. Consequently, $b(X\wr Q)=2$ if and only if $r(X)\ge D(Q)$. In particular, for $C_\ell$ of prime order $\ell$ acting regularly,
\begin{equation}\label{eq:cyc}
r(X\wr C_\ell)=\frac{r(X)^\ell-r(X)}{\ell}.
\end{equation}
\end{prop}

\begin{lem}\label{inherit}
Let $X$ be transitive on $\Delta$ with $b(X)=2$, let $Q\le S_m$ be transitive with $m\ge2$, and suppose $b(W)=2$ for $W=X\wr Q$. Suppose that, for every $\mathcal O\in\mathcal R(X)$ and every $\alpha,\beta\in\Delta$, there exist $\gamma,\delta\in\Delta$ such that $(\alpha,\gamma),(\gamma,\delta),(\delta,\beta)\in\mathcal O$. Then, for every $\mathcal P\in\mathcal R(W)$ and every $x,y\in\Delta^m$, there exist $u,v\in\Delta^m$ such that $(x,u),(u,v),(v,y)\in\mathcal P$. In particular, $\diam\Sigma(W)\le 3$.
\end{lem}

\begin{prop}\label{cprop}
Let $X$ be transitive on $\Delta$ with $b(X)=2$ and $r(X)\ge 2$, and let $\alpha,\beta\in\Delta$ be distinct points such that $\Sigma_X(\alpha)$ meets at most one regular $X_\beta$-orbit. Let $\ell$ be prime and let $W=X\wr C_\ell$ act on $\Delta^\ell$ in product action, with $C_\ell$ regular. Then $b(W)=2$, and $\alpha^\ell$ and $\beta^\ell$ are nonadjacent vertices of $\Sigma(W)$ with no common neighbour.
\end{prop}

\begin{lem}\label{criterion}
Suppose that $V_{\mathrm{reg}}\ne\varnothing$. Then:
\begin{enumerate}
\item[(i)] $b(G)=2$, and two distinct vectors $u,v\in V$ are adjacent in $\Sigma(G)$ if and only if $v-u\in V_{\mathrm{reg}}$.
\item[(ii)] Translation of an ordered base $(u,v)$ to $(\mathbf0,v-u)$ identifies the regular orbitals of $G$ with the regular $H$-orbits on $V$. In particular, $r(G)$ is the number of regular $H$-orbits on $V$.
\item[(iii)] Two vertices $u,v\in V$ have a common neighbour if and only if $v-u\in2V_{\mathrm{reg}}$. In particular, $G$ has the common neighbour property if and only if $2V_{\mathrm{reg}}=V$.
\item[(iv)] If $V_{\mathrm{reg}}\cup2V_{\mathrm{reg}}\ne V$ but $3V_{\mathrm{reg}}=V$, then $\diam\Sigma(G)=3$.
\end{enumerate}
\end{lem}

\begin{lem}\label{affineinput}
The group $L_0$ acts irreducibly on $U_0$, with precisely two regular orbits, namely $O_1$ and $O_2$. Moreover:
\begin{enumerate}
\item[{(i)}] $O_1+O_1=U_0$, and $|U_0\setminus(O_2+O_2)|=32$;
\item[{(ii)}] for $s,t\in\{1,2\}$, $\mathbf1\in O_s+O_t$ if and only if $s=t=1$.
\end{enumerate}
\end{lem}

\begin{equation}\label{eq:triple}
3O_1=3O_2=U_0.
\end{equation}

\begin{equation}\label{eq:tower}
U_{n+1}=U_n^3,\qquad
L_{n+1}=L_n\wr C_3,\qquad
G_{n+1}=U_{n+1}{:}L_{n+1}=G_n\wr C_3,
\end{equation}

\begin{thm}\label{thm:affine}
For every $n\ge0$, the group $G_n$ is a primitive affine group of degree $3^{9\cdot3^n}$ with $b(G_n)=r(G_n)=2$. If $n\ge1$, the vertices $\mathbf0$ and $\mathbf1$ of $\Sigma(G_n)$ are nonadjacent and have no common neighbour, and $\diam\Sigma(G_n)=3$. In particular, the groups $G_n$ for $n\ge1$ form an infinite family of affine counterexamples to Conjecture~\ref{bgconj}.
\end{thm}

\begin{proof}
Since $\dim_{\F_3}U_0=9$ and $\dim U_{n+1}=3\dim U_n$ by \eqref{eq:tower}, the degree of $G_n$ is $|U_n|=3^{9\cdot3^n}$. Lemma~\ref{affineinput} shows that $L_0$ is irreducible, so $G_0$ is primitive. Since $L_0\ne1$ and $L_{n+1}=L_n\wr C_3$, every $L_n$ is nontrivial, so $G_n=U_n{:}L_n$ is nonregular. Thus, if $G_n$ is primitive, then Lemma~\ref{paprim} shows that $G_{n+1}=G_n\wr C_3$ is primitive, since $C_3$ is transitive on three points. Moreover, $U_n$ is a regular elementary abelian normal subgroup of $G_n$, so $G_n$ is of affine type. Lemmas~\ref{criterion} and~\ref{affineinput} give $b(G_0)=r(G_0)=2$. Suppose that $b(G_n)=r(G_n)=2$. Since $D(C_3)=2$, Proposition~\ref{bookkeeping} gives $b(G_{n+1})=2$ and
\[
r(G_{n+1})=\frac{r(G_n)^3-r(G_n)}{3}=\frac{2^3-2}{3}=2.
\]

Now suppose that $\gamma$ is a common neighbour of $\mathbf0$ and $\mathbf1$ in $\Sigma(G_0)$. By Lemma~\ref{criterion}, there are $s,t\in\{1,2\}$ such that $\gamma\in O_s$ and $\mathbf1-\gamma\in O_t$. Since $\mathbf1=\gamma+(\mathbf1-\gamma)$, Lemma~\ref{affineinput}(ii) gives $s=t=1$. It follows that
\[
\Sigma_{G_0}(\mathbf0)\cap\Sigma_{G_0}(\mathbf1)\subseteq\mathbf1-O_1.
\]
Since scalar multiplication by $-1$ centralises $L_0$, the set $-O_1$ is again a single regular $L_0$-orbit, so $\mathbf1-O_1$ is a single regular $(G_0)_{\mathbf1}$-orbit. The vertices adjacent to $\mathbf1$ are precisely those lying in regular $(G_0)_{\mathbf1}$-orbits, so $\Sigma_{G_0}(\mathbf0)$ meets at most one such orbit and the pair $(\mathbf0,\mathbf1)$ satisfies the hypothesis of Proposition~\ref{cprop}. That proposition shows that the corresponding constant all-zero and all-ones vectors in $U_1$ are nonadjacent and have no common neighbour. These two properties pass to constant lifts from $G_n$ to $G_{n+1}$: nonadjacency follows from Lemma~\ref{basecrit}, and every coordinate of a common neighbour of the two lifts would be a common neighbour of $\mathbf0$ and $\mathbf1$ in $\Sigma(G_n)$. Thus the corresponding vectors in $U_n$ are nonadjacent and have no common neighbour for every $n\ge1$. Hence $\diam\Sigma(G_n)\ge3$ for $n\ge1$.

For the reverse inequality, let $\mathcal O_i$ be the regular orbital of $G_0$ corresponding to $O_i$, where $i\in\{1,2\}$. Given $x,y\in U_0$, \eqref{eq:triple} gives $a,b,c\in O_i$ such that $y-x=a+b+c$. Hence $(x,x+a),(x+a,x+a+b),(x+a+b,y)\in\mathcal O_i$. Repeated application of Lemma~\ref{inherit} shows that any two vertices of $\Sigma(G_n)$ are joined by a path of length at most three.
\end{proof}

\end{document}
